\subsection{Directorio de miembros y revocación durable}\label{sec:impl-miembros}

El directorio realiza el ciclo de vida de CU-02 y los roles y bajas de RF-02 y RF-03 (\cref{tab:cu-02,tab:rf-acceso}), aparte del directorio de personas de un expediente. El diseño, con la guarda del último Owner, está en la \cref{fig:cu-seq-alta}, y la revocación por generación, en la \cref{fig:impl-identidad-estado}. El alta directa por un Owner (\texttt{POST users}, \cref{tab:impl-http-identidad}) no es una invitación (\cref{sec:impl-invitaciones}): entrega el material MFA una sola vez y, si se pierde la respuesta, no se recupera.

\paragraph{Directorio.} Solo un Owner activo consulta el directorio y los selectores de asignados o disponibles, que muestran identificador, correo, rol, actividad y revisión, sin secretos ni generación. Los filtros preceden a una paginación de 1 a 100 filas por UUID, con cursor ligado a filtros y expediente y sin instantánea común entre páginas.

\paragraph{Cambios.} Rol y actividad cambian juntos contra la revisión esperada: un cambio real incrementa revisión y generación con su auditoría en una transacción, repetir valores vigentes no cambia nada y una revisión anterior da conflicto aunque coincidan. PostgreSQL impide dejar al despacho sin Owner activo, también ante cambios concurrentes. Desactivar conserva asignaciones y autoría; retirar una asignación es otra operación, permitida con el expediente cerrado. Reactivar exige contraseña y MFA, y un resultado incierto, volver a consultar la cuenta (\cref{sec:impl-qadra-conciliacion}).

\paragraph{Límites.} Una petición ya autorizada puede terminar tras la revocación, que no recupera bytes entregados. Restaurar PostgreSQL puede retroceder la generación: antes de reabrir se invalidan sesiones y desafíos del Redis correspondiente (\cref{anexo:miembros}). Ninguna primitiva coordina la restauración conjunta de PostgreSQL y Redis, descrita en la \cref{sec:impl-despliegue}, que los controladores de mantenimiento no automatizan del todo. La verificación está en la \cref{sec:pruebas-miembros}.

\Fuentes{\rutarepo{docs/members-api.md}; \rutarepo{docs/database-operations.md}; \rutarepo{docs/deployment-backups.md}; ADR 0043.}
